\section{Open-port comparisons and the deterministic distance scale}\label{sec:metric}

Write $\xi>0$ for the critical Weyl exponent of the reference metric, so that
$D_{h+b}=e^{\xi b}D_h$ for a constant $b$, and put
\[
 P_j=2^j,\qquad w_i=2^{-\xi i}\quad(i\in\Z).
\]
Initially the perimeter units are $P_j$. All graph paths in this section use
the original primal edges, including both types of blue diagonal in the
encoding. Their costs count whole edges. In denominator units tending to
infinity, splitting off partial edges at a fixed finite collection of
topological tests changes the total by a vanishing amount. An intersection
forced between two drawn primal walks is an intersection in the primal graph.

We work first on ordinary patches of the arbitrary-area, freely embedded
surface class. The reference metric on such a patch is the local critical
metric constructed in Section~\ref{sec:continuum}; its length structure,
Weyl rule, and ordinary topology ultimately come from
\citet{DGuniqueness,DGtopology}. No metric characterization theorem is applied
to a graph profile. The ingredients concerning graph observations are the
regularizations and conditional laws of Section~\ref{sec:kernels}, the actual
uniformization comparison in Section~\ref{sec:conformal}, and the field-kernel
comparisons established there. In particular, the source elimination and
smooth-field comparisons of Proposition~\ref{conf:source-elimination} and
Lemma~\ref{conf:field-kernel-locality} are already available.
Their area-one transfer will be used only for the local assertions obtained
below.  The optimal-comparison-constant and shortcut strategy adapts the
field-perturbation method of Ding--Gwynne
\citep[Sections~1.5 and~3--5]{DGuniqueness}.  Here the competing objects are
graph profiles, so the conditional-kernel arguments must also justify every
preservation of graph observations under a field change.

The maximal flag mesh tends to zero in the actual conformal coordinate.
Faces have vanishing diameter as well: all the flags incident to a face
share its center, so the diameter of the face star is at most twice the
largest incident flag diameter. Consequently an open Euclidean tube can be
followed topologically by primal paths. Winding paths can be closed by
forced transverse intersections with strict margins. These facts supply
topological witnesses; they supply no cost estimate for attaching a specified
lattice vertex to such a witness.

The proof has four stages. Shell circuits and traversal barriers first
convert high-probability annular tests into upper and lower path
comparisons. A first-failure argument then supplies growing denominator
units and record scales where an upper comparison is finite. For any such
extraction, the field modification argument proves equality of its optimal
upper and lower factors. Finally, uniform quantile retests propagate this
rigidity from record scales to every large perimeter scale. The outcome
is Proposition~\ref{metric:scale-bootstrap}; endpoint attachment remains
for Sections~\ref{sec:endpoints} and~\ref{sec:address}.

\subsection{The comparison statements and their test events}

\begin{definition}[Open-port comparison]\label{metric:comparisons}
Fix a joint extraction of regularized path observations, with deterministic
graph denominator units $T$ along that extraction.
An \emph{upper comparison by $C$} means the following almost-sure assertion
throughout the ordinary arbitrary-area class: every nonconstant reference
path $\pi$ of finite $D_h$-length, compactly contained in an ordinary patch, can be
followed with normalized graph cost at most $C$ times its reference length,
with every prescribed positive spatial, endpoint, and budget loss.
Explicitly, fix finitely many ordered open ports met by the reference
path, intervening open tubes containing its corresponding subpaths, and
\(\varepsilon>0\). Once these spatial margins and \(\varepsilon\) are
fixed, the extracted finite maps eventually have primal paths obeying
these instructions with cost divided by \(T\) at most
\(C\operatorname{len}_{D_h}(\pi)+\varepsilon\). The first and last
ports are open neighborhoods of the reference endpoints; their lattice
vertices are chosen as part of the witness.

A \emph{lower comparison by $c$} means that every sequence of tested graph
paths with a bounded normalized budget, confined to a compact ordinary test,
and with convergent endpoint locations has limiting inferior cost at least
$cD_h(x,y)$, where $x,y$ are the limiting endpoints. The observations include
marked subpaths and their costs, so that this assertion also applies to
every subpath with the same protections.
\end{definition}

All clearances are fixed before the comparison is made. Countable spatial
and budget margins suffice to formulate both assertions. For the upper
assertion, test the infimal reference length for each finite ordered
tube-and-port instruction, add strict slack, and then approximate an entire
path. The internal reference length data on buffered compact sets are
measurable from the fields. We use the optimal constants
\[
 C=\inf\{C':\text{upper comparison holds by }C'\},\qquad
 c=\sup\{c':\text{lower comparison holds by }c'\}.
\]
They are deterministic, since the defining assertions require probability
one throughout the specified measure class. Countable approximation retains
the weak inequalities at either optimum.

A positive lower comparison also gives an intrinsic lower bound in every
open enlargement of the compact test. Indeed, subdivide a bounded-cost path
into marked advances. The lower bound on subpaths and the ordinary topology
of $D_h$ prevent a fixed positive spatial advance from having vanishing cost.
At finer subdivisions, compactness therefore produces a continuous limiting
path, and sums of the subpath lower bounds control its reference length.
Equivalently, use finite chains of marked advances and let the mesh of these
chains tend to zero. Thus, if all the paths stay in a compact subset of an
open set $V$, their limiting inferior cost is at least $cD_h(x,y;V)$.

For every fixed integer offset $i$, the corresponding experiment at $j-i$,
with denominator $Tw_i$, has the same optimal constants against the prototype
reference metric. To see the units, adding $i\log 2$ to the real field
multiplies reference lengths by $w_i^{-1}$. The arbitrary-area input class is
equivalent under this global conversion. The product rule
\eqref{eq:original-15}, with its transformed cut data, gives the same spatial
readout rules in the converted parameters; the raw graph cost is unchanged.
Use simultaneous extractions for the countably many fixed offsets. Negative
$i$ gives the same assertion for fixed forward offsets. Probabilities at such
labels will always be evaluated in one common prototype distribution in its
converted units.

Fix a planar circular band in the prototype, with the ends conditioned to
lie outside a larger test disk. Let $A_j>0$ be a high quantile of the cost of
a primal circuit winding in a middle band. Choose this deterministic
quantile, with harmless approximation slack, so that for a fixed sufficiently
small $\eps_U>0$,
\begin{equation}\label{metric:quantile-tests}
 \Prob[\text{winding cost}\le 2A_j]\ge 1-2\eps_U,
 \qquad
 \Prob[\text{winding cost}\ge A_j/2]\ge \eps_U/2.
\end{equation}
Spatial errors tend to zero through the finite comparisons. Use a slightly
wider band for high-probability upper success and a narrower band when
retesting a limiting assertion. A winding circuit may be required to have
winding number of absolute value one; it contains a simple surrounding
subcircuit.

We also use a \emph{lower traversal test}: at a proposed threshold $t$, every
traversal between the two radii of a fixed circular band has cost at least
$t$, with probability at least $1-\eps_L$. Here $\eps_L>0$ is fixed and
sufficiently small. Nested bands provide slack in both directions when
passing to limiting observations. Arbitrary paths escaping outward contain
the required compactly confined subtraversals, so a lower test makes no
confinement assumption on the rest of such a path.

Strict limiting upper comparison implies upper success inside the
measurement band: take finitely many finite-length reference crossings with
transverse, strictly buffered intersections that stitch to a circuit in a
narrower band. Such paths have almost surely finite total length by the local
length-space topology. Strict limiting lower comparison gives lower success
by charging a positive advance in a compact ordinary subband, whose reference
crossing distance is bounded below in probability. These two observations
will give quantitative retests using the fixed marginal law of the prototype
reference metric.

\subsection{Chaining aligned shells}

\begin{lemma}[Shell chain rule]\label{metric:chain-rule}
Suppose that, along arbitrarily large intervals $[N,2N]$, a fixed positive
fraction of the offsets have sufficiently high-probability, buffered upper
circuits of cost at most $KTw_i$. Then the limiting observations have upper
comparison by $K/a_0$, for a deterministic $a_0>0$ depending only on the fixed
test geometries and probability tolerances. They also have arbitrarily small
simple surrounding primal circuits of arbitrarily small normalized cost
about every ordinary point.

If instead a fixed positive fraction have sufficiently high-probability
traversal barriers of cost at least $kTw_i$, then the lower comparison is at
least $k/b_0$, for a deterministic finite $b_0$. Both conclusions hold
together when both kinds of tests are imposed. If the available lower
factors tend to infinity with $N$, every positive advance costs unboundedly
in the chosen units.
\end{lemma}

\begin{proof}
Apply the aligned all-point shell estimate of
Proposition~\ref{conf:aligned-shells}. Include fixed reference-metric tests in the
annular buffers. For upper chaining, every path from the central hole, with
its fixed clearance, to the first possible circuit radius must pay at least
$a_0$ in prototype reference units; a strict subband gives this charge.
For lower chaining, take a round circle past the charged traversal band and
bound its internal $D$-diameter by $b_0$, using paths in a slightly larger
annular buffer. These tests can have arbitrarily high probability after
shrinking $a_0$ and enlarging $b_0$, by the ordinary topology and local length
structure. The fixed annular trial region can be made large enough to contain
all the required buffers. On an aligned shell at label $i$, these charges and
sizes are respectively $a_0w_i$ and $b_0w_i$.

With probability tending to one, the successful central holes give a finite
cover of each ordinary compact set. Choose an increasing sequence of $N$'s
along which these covers succeed almost surely. The covers are finite at
each $N$; fix one such cover before sending the discrete index to infinity.
All realization slack on the good couplings can be absorbed into $K$ and $k$.

For an upper comparison, begin at a point of the reference path in a good
hole and follow the path until its first exit from the filled disk bounded
by the associated simple circuit, or until the path ends. Every complete
step consumes at least $a_0w_i$ of reference length. At this fixed finite
mesh the minimum relevant $w_i$ is positive, so the procedure terminates.
Consecutive filled disks overlap. When two adjacent disks are nested,
compress their block to the containing representative, assigning to it the
reference steps of that block; all constituent disks stay inside that
representative. The remaining successive boundaries intersect. One can
therefore walk along the chain, traversing each boundary at most once.
The total normalized graph cost is at most
\[
 K\sum_{\text{complete steps}}w_i+O(K\max w_i)
 \le \frac K{a_0}\operatorname{len}_{D_h}(\pi)+O(K\max w_i).
\]
The compressed blocks still follow the original path in its prescribed
order, up to the uniformly shrinking diameters. The end locations have the
same approximation property. Passing through the finite discrete witnesses
and then sending the shell mesh to zero proves the upper comparison.
Applying the same cover to a single point supplies surrounding simple
circuits of arbitrarily small cost about every ordinary point.

For a lower comparison, follow a tested graph path from a hole to its first
exit on the indicated round circle, and repeat. Each complete step costs at
least $kw_i$ in normalized graph units. A fixed bounded total budget bounds
the number of these steps at a fixed mesh; marked-subpath extraction handles
any variation in the witnesses. The limiting circles, after the same
containment reduction, form a chain connecting arbitrarily close to both
endpoints. Its reference cost is at most
\[
 b_0\sum_{\text{complete steps}}w_i+O(b_0\max w_i).
\]
The last incomplete step is included in the error. Reference continuity
removes the remaining endpoint errors. This uses no bound on the diameter
of a shell's interior in that shell's units. It follows that the lower
comparison is at least $k/b_0$. The same argument with arbitrarily large $k$
gives infinite lower comparison. In bounded raw graph units, infinite lower
comparison already follows from vanishing primal-edge mesh.
\end{proof}

The probability tolerances in Lemma~\ref{metric:chain-rule} can be fixed
uniformly for the basic tests, independently of the extraction and of the
offset. Fix, for example, the required fraction at $1/10$, with a fixed
reduction to leave indexing margins. Only a fixed collection of band
geometries is involved. If no probabilities sufficiently close to one worked
uniformly, put the countably many contrary extractions on common prototype
inputs. Theorem~\ref{kern:finite-vector} applies to their
\emph{entire} kernel vector.
Truncate its minimum positive weight once for that vector. The density and
bounded harmonic changes used in the shell estimate then turn the proposed
high prototype probabilities into adequate simultaneous-candidate success
probabilities for one of the contrary extractions, a contradiction. The
reference regularity cutoffs depend only on the common field and can also
be fixed uniformly. This argument requires one finite index for the whole
countable vector, rather than separate finite supports for its coordinates.

Quantiles and traversal tests pass from finite comparison probabilities to
limiting rules by further joint extraction with intermediate open bands and
covers. For winding witnesses retain ordered subpaths and their robust
intersections. As with the extremal-length barrier
\eqref{eq:original-16}, no unregularized chart cost is assumed to be a
continuous function of its chart data.

\begin{lemma}[Initial growth and record scales]\label{metric:records}
For every $\delta\in(0,\xi)$ the scales $A_j$ have exponential lower growth
with exponent $\xi-\delta$ in base two. There is a sequence of record starts
$j\to\infty$ along which
\begin{equation}\label{metric:record-ratios}
 \limsup\frac{A_{j-i}}{A_j}\le w_i
 \qquad\text{for every fixed positive integer }i.
\end{equation}
\end{lemma}

\begin{proof}
Starting at $j_0$, test lower traversal at the raw threshold
$2^{(\xi-\delta)(j-j_0)}$. As $j_0\to\infty$, vanishing mesh supplies every
fixed finite initial run of successful tests. If arbitrarily large starts
had a later first failure, extract at those failures with $T$ equal to the
proposed threshold. Their successful preceding runs have diverging length,
and every fixed preceding offset has a high barrier of size
\[
 \text{constant}\cdot Tw_i2^{\delta i}.
\]
For $i\in[N,2N]$ this gives lower factors tending to infinity as $N\to\infty$.
Lemma~\ref{metric:chain-rule} gives infinite lower comparison and hence
strictly re-passes the allegedly failed test. Thus all sufficiently large
starts have no first exit.

Fix such a start and extract at arbitrary $j\to\infty$ in its proposed
threshold units. The same argument gives infinite lower comparison. In
particular, winding circuits cannot have bounded cost in those units with
positive limiting probability. The upper probability in
\eqref{metric:quantile-tests} proves the asserted lower growth of $A_j$.

To obtain records, choose exponents tending to $\xi$ from below. For an
exponent $\xi-\delta$, the preceding lower growth at the slightly larger
exponent $\xi-\delta/2$ shows that
$A_j2^{-(\xi-\delta)j}$ has arbitrarily late records. At each record,
$A_{j-i}/A_j\le2^{-(\xi-\delta)i}$ for preceding indices. Let $\delta\downarrow0$
and the record index tend to infinity diagonally. This gives
\eqref{metric:record-ratios} simultaneously at all fixed positive offsets.
\end{proof}

\subsection{Certificates and a supply of shortcuts}

For the rigidity argument, fix an extraction with $T\to\infty$ for which
upper chaining gives a finite upper constant and for which the optimal
constant satisfies $C>0$. The same comparisons and surrounding-circuit
properties hold at every fixed offset in $Tw_i$ units.

For an ordinary open set \(U\), a path certificate of budget \(t\)
from \(x\) to \(y\) records limiting marked-path observations whose
range lies in some compact \(K\Subset U\), whose endpoints tend to
\(x,y\), and whose normalized cost is at most \(t\). Its meaning is
through the buffered observations: for every positive endpoint and budget
loss and every open enlargement of \(K\) in \(U\), actual primal
witnesses eventually realize those relaxed requirements in the extraction.
Intermediate marked subpaths and their budgets are retained as well.

Let \(F_U(x,y)\) be the infimum of these budgets, taking the diagonal
closure as the losses decrease to zero and then the infimum over
\(K\Subset U\); set \(F_U(x,x)=0\). Countable compact protections,
open enlargements, and rational losses suffice. Thus a certificate is an
assertion about actual finite-map witnesses with open margins. Its
endpoints do not designate particular lattice vertices.

\begin{lemma}[Joining certificates]\label{metric:certificate-joining}
For points in the same component of $U$,
\begin{equation}\label{metric:certificate-upper}
 F_U(x,y)\le C D_h(x,y;U).
\end{equation}
Moreover $F_U$ obeys the triangle inequality and is continuous on local
endpoint compact sets within that component. The same joining rule applies
with intermediate ordered-port constraints having strict margins.
\end{lemma}

\begin{proof}
Follow a reference path and then take its infimal length to obtain
\eqref{metric:certificate-upper}. To join two nonconstant certificates at a
common endpoint, choose a surrounding simple primal circuit in $U$ of
arbitrarily small diameter and cost, supplied by
Lemma~\ref{metric:chain-rule}. Cut the two certificates at this circuit and
join along it. At each fixed circuit and protection, use simultaneous path
witnesses with slack. Even a separately obtained circuit forces the required
intersections, by its central-hole clearance. Compact confinement is
preserved. Let its cost and diameter tend to zero to obtain the triangle
inequality. The upper bound and the reference length topology then give
continuity in each endpoint, and hence jointly on local compact sets.
Ordered ports survive the same construction when their margins are fixed.
\end{proof}

The lower optimum satisfies $0\le c\le C$: follow a short nontrivial reference
geodesic in a small ordinary patch, or paths arbitrarily close to minimizing,
and apply the lower bound to the upper witnesses with their endpoint and
range approximations.

\begin{lemma}[Strict rectangle shortcuts]\label{metric:shortcut-supply}
Suppose that $c<C$, and fix $c_1\in(c,C)$. For all sufficiently large $N$, at
least a fixed positive fraction, which may be taken as one half, of the
labels $i\in[N,2N]$ have probability at least $p>0$ of the following event.
In $Tw_i$ units there is a path certificate in a rectangle $R$, in a compact
middle subband of an ordinary annular test, joining points $q_0,q_1$ on
opposite faces, with cost at most
\[
 c_1D_h(q_0,q_1;\Omega).
\]
Here $R\Subset\Omega$, with uniform clearance, where $\Omega$ is a fixed
larger ordinary annular buffer. Rectangles may have a fixed small crossed
width, bounded transverse extent, variable centers, and either coordinate
axis as crossing direction. The certificate uses a closed rectangle
envelope and arbitrary open spatial, endpoint, and cost losses.

The same conclusion holds when the tested labels are $i+m$ for any fixed
integer $m$, with a slightly smaller fixed fraction.
\end{lemma}

\begin{proof}
Otherwise, on a fixed positive fraction of labels along arbitrarily large
intervals, absence of all these shortcuts has arbitrarily high probability.
Use this absence as a shell test. Include the upper circle-diameter cutoff
$b_0$ and a uniform positive lower bound $a_1$ for reference internal distances
between opposite faces of the candidate rectangles. These are compact
parameter tests and can be imposed with countable spatial refinements.

In every complete lower-chain step, mark a rectangular subtraversal as
follows. In the middle band start a small centered square, follow the path
to its first face exit, and retain the segment after its last crossing of
the parallel midline. The resulting half-square is a permitted $R$.
By compact marked extraction, absence of a shortcut forces this portion
to pay more than $c_1D_h(q_0,q_1;\Omega)$, with a slightly diminished strict
threshold if necessary. Since $D_h(q_0,q_1;\Omega)$ dominates the unrestricted
distance, the excess over the original \(cD_h\) endpoint charge is
at least \((c_1-c)a_1w_i\) before the strict-loss adjustment. Fix, for
example, \(0<\delta_1<(c_1-c)a_1\), and absorb that adjustment into the
remaining gap. The extra charge is then at least \(\delta_1w_i\). In particular, a bounded budget
bounds the number of complete steps at a fixed mesh even when $c=0$.

Retain the original lower comparison on all intervening portions. Summing
their endpoint distances and using the triangle inequality gives a total
charge at least $cD_h(x,y)+\delta_1\sum w_i$, up to the last incomplete-step
error. The circle chain gives
$D_h(x,y)\le b_0\sum w_i+o(1)$. Hence the lower factor improves by a fixed
positive amount, contradicting the definition of $c$. Fixed label shifts
discard only boundedly many labels and are handled by the same unit
conversion; the argument is unchanged. All constants in the strict gap
$C-c_1$ are fixed before any later near-saturation tolerance.
\end{proof}

\subsection{A protected shortcut forced by a smooth field modification}

We continue under the contrary assumption $c<C$, with $c_1$ fixed as in
Lemma~\ref{metric:shortcut-supply}. We first construct one local modification
and then count its possible effects on a near-saturating reference geodesic.
Each successful modification will force a reference geodesic to use a
short interval where a preserved graph certificate saves a fixed amount
against upper factor \(C\). The subsequent visit count will compare the
weight of these forced savings with the small total saving allowed by a
near-saturating pair of endpoints. The geometry, inner scale, and
change-of-law bounds are all chosen before that near-saturation tolerance.

The modification concerns the reference metric through its Weyl rule and
the graph observations through their field kernels. It does not assert a
Weyl rule for graph costs in a region where the added function varies.

\paragraph{Outer and inner units.}
An outer shell has label $i$, reference length unit $w=w_i$, and a small
central hole. Center its coordinate and divide by its spatial radius, pull
back the real field with charge two, and add $i\log2$. Take its attachment
circle to have radius $2$ in this coordinate. Inner shortcuts will use label
$i+m$, where $m$ is large but fixed, and their length unit is
\begin{equation}\label{metric:inner-unit}
 u=w_{i+m}=w\,2^{-\xi m}.
\end{equation}
The outer tests require a reference charge at least $aw$ from the central
hole to the attachment circle on either side of a visit, and internal
diameter at most $K_0w$ for that circle in a slightly larger annular buffer.
Include the ordinary length and continuity cutoffs on compact subannuli
used below. The constants can be enlarged, or $a>0$ decreased, when imposing
these high-probability cutoffs. All shifts will be supported strictly inside
the attachment circle and outside the central omitted disk.

\paragraph{Robust inner shortcuts.}
Place small compartments with their initial test buffers in a middle open
band of the outer shell. At a temporarily shifted real height $h+b$, retain
the shortcut event of Lemma~\ref{metric:shortcut-supply} together with
\begin{equation}\label{metric:shortcut-size}
 d=D_{h+b}(q_0,q_1;\Omega)\in[a'u,A'u],
 \qquad \text{certificate cost}\le c_1d,
\end{equation}
where $0<a'<A'<\infty$. These are actual length and base denominator units
after placing the prototype in the compartment; symbols such as $\Omega$
and $R$ denote their placed images when used in those units. Loose cutoffs
for the common prototype metric preserve a fixed positive part of the
shortcut probability.

Choose
\[
 e=e_*u,
\]
with $e_*>0$ sufficiently small compared with $(C-c_1)a'$, with $a'$, and
with the error allowed in joins charged at factor $C$. Impose the following
finite geometric and metric refinements, still retaining a fixed positive
part of the probability in Lemma~\ref{metric:shortcut-supply}.
There are disjoint larger terminal caps at $q_0,q_1$ whose connection errors
to the corresponding endpoint, within $\Omega$, are at most $e$ for
$D_{h+b}$. Points sufficiently closer to $q_j$ can be joined to it inside a
much smaller safe subcap within the same budget. Use finite position bins
for the rectangle and its endpoints, and nested thin enlargements of the
rectangle. One such enlargement is a preservation neighborhood $V$ with
strict clearance about $R$. There is a path in $V$ from $q_0$ to $q_1$ of
$D_{h+b}$-length at most $K_1u$. Terminal stops can be placed just outside a
larger chamber enlargement, near its corresponding faces.

Here is the order for choosing these refinements. On compact subsets of the
prototype $\Omega$, first use continuity to choose a small larger-cap size
with high probability of the required error bound. Choose a closer-point
scale for joins inside the safe subcaps, using continuity and a positive
charge for leaving the smaller cap radius. Choose stop distances, the thin
margins of $V$ and of the chamber, and then position bins at still finer
scales. A finite length cutoff inside $V$ is obtained from connected inner
compact corridors with room around those bins. All restrictions are chosen
using the common marginal reference metric law. No regularity of graph
endpoints, and no uniform cost estimate for a passage of vanishing width,
is used here.

Also require a whole bounded-side $O$ disk enclosing the inner test compact,
with that disk contained in a larger fixed ball in inner coordinates. By
the post-conformal field-kernel locality of
Lemma~\ref{conf:field-kernel-locality},
the shortcut success function is then determined by this field buffer.
Choose the buffer large using the common marginal law of $J$. End marks
and the central choice can be redrawn outside the enclosing disk. A law
with ends even outside the field buffer has the same supported unmarked
field-kernel possibilities: the admissible end choices have positive
conditional mass, and source and central-choice elimination has already
been proved. Truncating the resulting densities preserves a uniform
positive probability. Thus this local determination event and the refined
shortcut have probability bounded below in ordinary actual-height GFF
classes on the larger ball, strictly inside the observation region.

\paragraph{Amplification within an outer shell.}
Pack many similar compartments, with disjoint Dirichlet buffers, into the
outer annulus. In the normalized outer coordinate a compartment of
similarity scale $s$ uses, for its inner field comparison, the further
pullback and addition
\[
 2\log s+m\log2+b.
\]
Choose $b$ from rounded constants that compensate the local coarse real
height; the common term $-m\log2$ is included in this choice. After the
number and sizes of the compartments have been fixed, truncate their
compensating constants to finite lists at arbitrarily small probability
cost. Given the fields, draw private auxiliary shortcut profiles from the
field kernels at the corresponding globally shifted inputs. The fixed
offset comparisons justify their use at label $i+m$.

For amplification, condition outside the compartment buffers in the
ordinary Gaussian experiment. With high probability most compartments
have bounded real harmonic differences in inner scale after centering,
and bounded imaginary harmonic data there modulo its period. The marginal
Gaussian harmonic bounds are uniform in scale; they are used before the
number of compartments is chosen. Absolute real constant ranges can be
truncated afterwards. In each of these good compartments, independent
Dirichlet interiors with the bounded harmonic corrections dominate a
fixed part of the positive-probability inner event, by the fixed-class
absolute-continuity comparison and the Cameron--Martin inequality with
bounded-norm cutoffs. The rounded compensating shift removes dependence
on the absolute coarse real height and on $m$. The lower bound for this
conditional success probability is therefore independent of the packing
size and of $m$.

Count only successes on the local determination event. Conditional
independence of the Dirichlet interiors and of the private proposals then
amplifies this positive probability to a probability arbitrarily close to
one. Any initial enclosing density change is handled by truncation. The
geometry, $K_1$, and $a'$ can be chosen independently of $m$, except for
the compensating constant lists and the label $i+m$. Choose $m$ now so
large that
\begin{equation}\label{metric:short-budget}
 B_0=(K_1+12e_*)u\ll aw.
\end{equation}
Use only outer labels having the inner probability supplied by
Lemma~\ref{metric:shortcut-supply}; they still form a fixed positive
fraction. The outer, packing, and additional recipe regularity failures
are chosen sufficiently small for the aligned-shell theorem, with this
fixed fraction and a fixed encompassing annulus. The proposal success
probabilities on their determination events are functions of the local
fields. Proposals may also be sampled privately over shells. Consequently
the aligning field changes use ordinary local equivalences and the affine
and offset kernel comparisons, without a global order selection. The trial
count can be increased to meet any prescribed fixed tolerance.

\begin{lemma}[Finite modification recipes]\label{metric:forcing-recipe}
For the successful shells just constructed there is a fixed finite list of
smooth field modifications, in normalized outer coordinates, with the
following property. Let $W$ be an ordinary open patch containing the shell
with room. Suppose that $p_0,p_1\in W$ lie outside its enlarged disk and
that a $D_h(\cdot,\cdot;W)$-geodesic of length $L$ passes through the central
hole. A recipe $f$ from the list preserves the selected inner certificate,
as a certificate for $h'=h+f$, and preserves all old profile readings
compactly outside a strictly enlarged shell disk.

Every minimizing $D_{h'}(\cdot,\cdot;W)$-geodesic, when it exists, has a
marked interval of length at most $B_0$ with an $F'_U$ saving at least
$\delta_*u$ against upper factor $C$, provided the geodesic is compactly
contained in an upper-comparison domain $U$ containing this shell with
room. Here $\delta_*>0$ is fixed, and at least $a'u/2$ of the marked
interval lies inside the enlarged shell disk. The finite list and all
constants are independent of the later near-saturation event and of its
probability.
\end{lemma}

\begin{proof}
Let $l_j$ be the infimal cost in $W$ to first hit the attachment circle from
$p_j$, and choose such hits $x_j$ with errors less than $e$. Since the old
geodesic enters the central hole,
\begin{equation}\label{metric:old-gap}
 l_0+l_1+2aw\le L.
\end{equation}
The points $x_0,x_1$ are separated at a fixed positive outer Euclidean
scale: otherwise the outer continuity cutoff would join them in the
annular buffer at cost less than $aw$, contradicting
\eqref{metric:old-gap}. Put the $x_j$ in finite position bins, with larger
disjoint outer caps having old connection error at most $e$ to $x_j$.
Choose outer stops just inside circle $2$, with safe smaller-cap joins of
cost at most $e$. These outer cutoffs are chosen after $m$, in the old
field, in the same order as the inner caps. Only finitely many descriptions
are needed.

We specify the geometry of the recipe. Take an effective open union
\[
 O\cup T_0\cup S\cup T_1.
\]
Here $O$ is the exterior of a circle of radius slightly less than $2$;
$S\Subset\Omega$ is a chamber about $V$ in the selected compartment; and
$T_0,T_1$ are disjoint tubes connecting the corresponding outer and
terminal caps. The only overlaps are between $O$ and a tube, or $S$ and
a tube, inside the relevant larger caps and strictly inside circle $2$.
Avoid the omitted central disk. Choose a smaller effective union with
closure inset in the larger one, still containing the attachment circle
and exterior, $V$, the compact safety zones for cap joins, and tube cores
connecting the chosen stops. A closed layer of positive thickness
separates this smaller union from exit out of the larger one.

The core stops lie outside $\overline O\cup\overline S$. Their joins to
$x_j$ and $q_j$ stay in the safe subcaps. The parts of the tubes to be
lowered are compactly supported away from $\overline O\cup\overline S$.
Choose a thick transverse marker strip in each tube between its two
overlaps, away from the end regions.

For completeness, the required disjoint layout can be chosen with strict
margins. Draw two arcs from just inside the outer circle to just outside
opposite faces of the chamber enlargement, avoiding the central omitted
disk and prescribing short end approaches. After cutting along the first
arc connecting the two distinct boundary components, the remaining
available region is connected for the second arc. Thin sausages about the
arcs give the tubes, with disjoint enlarged end-cap neighborhoods. The
threshold defining $O$ can be arbitrarily close to $2$, so tube ends
overlap $O$ but stop strictly inside circle $2$. Include in the smaller
tubes every part of a safe outer join not already in the smaller exterior.
Similarly, a thin chamber enlargement permits terminal stops outside it
within the closer-point scale, with safe joins in the inner effective
union and overlaps within the larger terminal caps. Choose position bins
only after these scales. All radii and placements have clearance from
support obstructions. Thus the separating layer and marker strips do not
interfere with safe joins. This construction asks for no approach to an
unspecified side of a graph certificate beyond its protected rectangle
envelope.

Figure~\ref{fig:protected-shortcut} records the separated regions just
constructed. The diagram concerns the support and plateaux of the field
modification; metric costs are imposed in the next step.

\begin{figure}[htbp]
\centering
\begin{tikzpicture}[x=2.05cm,y=2.05cm,font=\small,
  line cap=round,line join=round,
  guide/.style={black!60,thin},
  refpath/.style={black!70,densely dashed,thick,
    -{Stealth[length=1.6mm]}},
  shortcutcap/.style={draw=blue!45!black,densely dotted,thin},
  note/.style={align=center,font=\small}]
\definecolor{shortcutwall}{RGB}{241,207,197}
\definecolor{shortcutmarker}{RGB}{204,119,79}
\definecolor{shortcutcore}{RGB}{25,125,105}
\definecolor{shortcutcert}{RGB}{110,57,144}

\draw[black!65,densely dashed] (0,0) circle (2);
\draw[black!45,densely dotted] (0,0) circle (1.80);

\fill[shortcutwall,even odd rule]
  (0,0) circle (1.88) (0,0) circle (1.82);
\draw[shortcutwall,line width=19pt,rounded corners=9pt]
  (-1.78,.25) -- (-1.50,.25) -- (-1.17,1.10) -- (-.57,1.10);
\draw[shortcutwall,line width=19pt,rounded corners=9pt]
  (1.78,.25) -- (1.50,.25) -- (1.17,1.10) -- (.57,1.10);
\fill[shortcutwall,rounded corners=2pt]
  (-.70,.76) rectangle (.70,1.44);
\draw[white,line width=12pt,rounded corners=9pt]
  (-1.82,.25) -- (-1.50,.25) -- (-1.17,1.10) -- (-.57,1.10);
\draw[white,line width=12pt,rounded corners=9pt]
  (1.82,.25) -- (1.50,.25) -- (1.17,1.10) -- (.57,1.10);
\fill[white,rounded corners=1pt]
  (-.63,.83) rectangle (.63,1.37);

\fill[white] (0,0) circle (.43);
\draw[black!55,densely dashed] (0,0) circle (.43);
\node[note,font=\footnotesize] at (0,0)
  {omitted\\central disk\\$f=0$};

\draw[shortcutcap] (-1.84,.25) ellipse [x radius=.25,y radius=.23];
\draw[shortcutcap] (1.84,.25) ellipse [x radius=.25,y radius=.23];
\draw[shortcutcap] (-.64,1.10) ellipse [x radius=.25,y radius=.23];
\draw[shortcutcap] (.64,1.10) ellipse [x radius=.25,y radius=.23];

\begin{scope}[shift={(-1.35,.64)},rotate=69]
  \fill[shortcutmarker] (-.055,-.15) rectangle (.055,.15);
\end{scope}
\begin{scope}[shift={(1.35,.64)},rotate=-69]
  \fill[shortcutmarker] (-.055,-.15) rectangle (.055,.15);
\end{scope}

\draw[shortcutcore,line width=3.2pt,rounded corners=9pt]
  (-1.63,.25) -- (-1.50,.25) -- (-1.17,1.10) -- (-.84,1.10);
\draw[shortcutcore,line width=3.2pt,rounded corners=9pt]
  (1.63,.25) -- (1.50,.25) -- (1.17,1.10) -- (.84,1.10);

\filldraw[fill=violet!6,draw=shortcutcert!70,thin]
  (-.56,.94) rectangle (.56,1.28);
\draw[shortcutcert!60,densely dotted,thin]
  (-.45,1.015) rectangle (.45,1.17);
\node[font=\footnotesize] at (0,1.225) {$V:\ f=b$};
\draw[shortcutcert,very thick]
  (-.45,1.10) -- (-.30,1.065) -- (-.13,1.135)
  -- (.04,1.07) -- (.23,1.13) -- (.45,1.10);
\fill[shortcutcert] (-.45,1.10) circle (1.2pt);
\fill[shortcutcert] (.45,1.10) circle (1.2pt);
\node[below,font=\scriptsize] at (-.45,1.095) {$q_0$};
\node[below,font=\scriptsize] at (.45,1.095) {$q_1$};
\node[right,font=\scriptsize,text=shortcutcert] at (.45,1.17) {$R$};

\draw[refpath,rounded corners=9pt]
  (-2.25,.25) -- (-1.50,.25) -- (-1.17,1.10) -- (-.69,1.10);
\draw[refpath,rounded corners=2pt]
  (-.69,1.10) -- (-.60,1.10) -- (-.60,1.335)
  -- (.60,1.335) -- (.60,1.10) -- (.69,1.10);
\draw[refpath,rounded corners=9pt]
  (.69,1.10) -- (1.17,1.10) -- (1.50,.25) -- (2.25,.25);

\node[above] at (0,1.49) {$S:\ f\ge b$};
\node[above left] at (-1.17,1.16) {$T_0$};
\node[above right] at (1.17,1.16) {$T_1$};
\node[note] at (0,2.30) {$f=0$ near and outside $|z|=2$};
\node[note,text width=2.7cm] at (2.78,1.39)
  {old exterior\\profile readings\\retained};
\node[note] at (2.40,-.65) {$O:\ f\ge0$};
\draw[guide] (2.13,-.65) -- (1.86,-.57);
\node[note,text width=2.2cm] at (-2.64,-.60)
  {high\\separating\\layer};
\draw[guide] (-2.20,-.60) -- (-1.75,-.60);
\node[below] at (0,-2.02) {attachment circle $|z|=2$};
\node[note,font=\footnotesize] at (0,-1.35)
  {$\partial O$: radius strictly below $2$};
\draw[guide] (0,-1.53) -- (0,-1.80);

\draw[shortcutcore,line width=3.2pt] (-2.70,-2.48) -- (-2.28,-2.48);
\node[right,font=\footnotesize] at (-2.19,-2.48) {lowered reference core};
\fill[shortcutmarker] (.55,-2.53) rectangle (.95,-2.43);
\node[right,font=\footnotesize] at (1.04,-2.48) {transverse marker};
\draw[shortcutcert,very thick] (-2.70,-2.78) -- (-2.28,-2.78);
\node[right,font=\footnotesize] at (-2.19,-2.78) {graph certificate in $V$};
\draw[refpath] (.55,-2.78) -- (.95,-2.78);
\node[right,font=\footnotesize] at (1.04,-2.78) {reference transit};
\end{tikzpicture}

\caption{Regions used by one field modification, shown schematically in
outer-shell coordinates. The graph certificate is preserved in \(V\),
where \(f=b\), and old exterior profile readings are preserved beyond
the shell, where \(f=0\). The two disjoint tubes and chamber constrain
reference-metric geodesics: high layers obstruct exits, while lowered
cores and safe cap joins provide access to the terminal caps. The central disk is outside the support.
Shapes and widths represent separation and incidence, not metric distances.}
\label{fig:protected-shortcut}
\end{figure}

\Needspace{6\baselineskip}
Choose a smooth baseline bounded above by a large $H_0$, nonnegative on
$O$, and at least $b$ on $S$. It equals $b$ on the preservation
neighborhood, equals $0$ on the outer safety zones, and equals $b$ on the
terminal safety zones; set it to $H_0$ on the forbidden layer and marker
strips. The clearances just established allow these plateaux. First impose
loose positive old-metric lower cutoffs for crossing the fixed forbidden
gaps and marker strips in $w$ units, and then choose $H_0$ sufficiently
large. A crossing at the unreduced wall or marker height now costs more
than $B_0$.

Subtract $J_0\chi$ from the baseline, where $\chi\ge0$ is smooth and
supported compactly in the smaller tubes, away from the wall layer and
$\overline O\cup\overline S$. In each tube its positivity set is connected
and contains a longitudinal core and its two stops. Choose $J_0$ large.
The following two properties can be imposed by finitely many further
high-probability old-metric cutoffs:
\begin{enumerate}
\item Each core connects its stops at modified cost at most $e$.
\item Every transit from outer to terminal overlap within the tube, or
in the reverse direction, of modified cost at most $B_0$ visits a point
with a modified join of cost at most $2e$ to the corresponding $q_j$,
compactly inside the shell enlargement.
\end{enumerate}
To prove the second property, a transit must meet the core-lowering region
within the marker; otherwise the marker charge exceeds $B_0$. The closure of the possible encounters in a marker is compact.
Choose a finite Euclidean net for it with centers in \(\{\chi>0\}\);
these centers exist because the encounter set lies in that positivity
set. By a further old internal-modulus cutoff, every encounter joins a
net center for cost at most \(e/2\), even at the upper height \(H_0\).
Join the finitely many net centers and stops by paths compactly in the
connected open positivity set. Each chosen path has a positive minimum
of \(\chi\). Impose a deterministic upper cutoff on these finitely
many old lengths and a deterministic positive lower cutoff on these
minima, at arbitrarily small additional probability cost. Then choose
one deterministic finite \(J_0\) making all these connections cheap. No positive lower bound
for \(\chi\) at every possible encounter is needed. Finish with the terminal safe-cap join. This proves the
claim without requiring metric regularity at the exact boundary of a
cheap plateau.

The resulting $f$ is $b$ near the preserved rectangle and is zero near and
outside circle $2$. All amplitudes and cutoff functions, in outer
coordinates, range over fixed finite lists. Coarse-height compensation and
the inner proposals decide only which entries of those lists are useful.
For each recipe also impose a bounded Gaussian change-of-law likelihood
test on its support. Since the list is finite, these tests have
arbitrarily high marginal probability after their cutoffs are enlarged.
They are imposed before any later geodesic event is selected.

Use the conditional kernel comparison of
Lemma~\ref{conf:field-kernel-locality} to
retain the selected certificate from $h+b$ inside the preservation region
and the old profile readings outside a strictly enlarged shell disk in the
modified field $h'=h+f$. The old reference prefixes to $x_j$, safe joins,
cheap cores, and the preserved connection in $V$ give a modified reference
path of length at most
\begin{equation}\label{metric:modified-upper}
 l_0+l_1+B_0<L.
\end{equation}
Here $12e$ accommodates all joining errors in addition to $K_1u$.

Consider a minimizing modified geodesic. It must enter circle $2$ to beat
$L$, since the metric outside that circle is unchanged. Its portion
between first and last visits to the circle has length at most $B_0$:
the unchanged exterior prefixes cost at least $l_0$ and $l_1$.
That portion cannot leave the effective union, because it would cross
the forbidden layer at its unchanged high height and pay more than its
budget. Its first effective tube use from $O$ is on side $0$, and its
last return to $O$ is on side $1$. For if a prefix from $p_0$ in $O$
reached the wrong mouth, it would cost at most $l_0+B_0$; since $f\ge0$
on $O$, the same prefix costs no more in the old field. Complete it by
the old cap join and the old prefix from $x_1$ to $p_1$. This contradicts
\eqref{metric:old-gap}, by \eqref{metric:short-budget}. The reversed
argument handles the last mouth. A transition between the two mouths
through $O$ during the middle portion is excluded by the same local
budget and cap-joining estimate.

Take a finite itinerary subordinate to the effective open cover, switching
only in overlaps. Some excursion from $O$ returns through a different tube
than the departure tube, since remaining in $O$ is impossible and the
first and last types differ. On this excursion choose a transition between
distinct terminal-cap types. This supplies a full approach through one
tube, a full departure through the other, and a portion assigned wholly
to $S$ between opposite terminal caps; spur returns do not alter this
choice. Since $f\ge b$ in $S$, the chamber portion alone costs at least
$d-2e$. The marker accesses on the approach and departure mark an interval
containing it, of length at most $B_0$, whose endpoints have modified joins
of cost at most $2e$ to $q_0,q_1$.

When this geodesic is compactly contained in $U$, the triangle inequality
of Lemma~\ref{metric:certificate-joining}, upper comparison in the modified
law, and the preserved certificate give $F'_U$ cost at most
\[
 c_1d+4Ce
\]
between those interval endpoints. The saving against upper factor $C$
is at least
\[
 C(d-2e)-(c_1d+4Ce)
 \ge \bigl((C-c_1)a'-6Ce_*\bigr)u.
\]
Choose $e_*$ so that the right-hand coefficient is some fixed
$\delta_*>0$ and $d-2e\ge a'u/2$. The chamber portion is inside the
enlarged disk, proving also the stated in-disk length bound. Upper
comparison holds in the modified law by full-field law domination.
Only the reference metric was compared by height inequalities in the
transition regions; the graph certificate was used solely where its
kernel is preserved.
\end{proof}

\subsection{Rigidity by counting modified geodesic visits}

\begin{proposition}[Equality of local comparison factors]\label{metric:rigidity}
In an extraction with diverging denominator units for which upper chaining
gives a finite positive upper optimum $C$, the lower optimum equals $C$.
The common factor is deterministic. This conclusion holds simultaneously
in the fixed-offset experiments with their converted units.
\end{proposition}

\begin{proof}
Suppose $c<C$, and retain the shortcuts, probability tolerances, inner
offset $m$, and finite recipes constructed above. All of these choices are
fixed before choosing the small near-saturation tolerance $\eta$ below.

\smallskip\noindent\emph{A near-saturating compact experiment.}\par
For arbitrarily small $\eta>0$, optimality of $C$ supplies deterministic
ordinary patches $U\Subset W$, a compact protection $K\Subset U$, fixed
anchors $p_0,p_1\in K$, and a positive-probability event $E$ on which
\begin{equation}\label{metric:saturation-event}
 \begin{gathered}
 A<L<A_1,\qquad A>0,\quad A<A_1\le2A,\qquad
 L=D_h(p_0,p_1;W),\\
 \min_{j=0,1}D_h(p_j,W\setminus K;W)>s_0L,\qquad
 F_U(p_0,p_1)>(C-\eta)L,
 \end{gathered}
\end{equation}
where $s_0>1$ has fixed slack. To justify this localization, if there were
no such witnesses, subdivide any reference path and its open following
instructions into sufficiently short subarcs in small ordinary domains
having that exit safety. Compactness, the length topology, and positivity
of exit costs permit arbitrarily fine subdivision in reference length.
Apply the resulting improved $F_U$ bound there, using endpoint continuity
to reduce to countably many anchor tests, and concatenate with
shrinking-cost circuit joins inside the prescribed tubes. This would
improve upper followability strictly below $C$, a contradiction. Restrict
the length to one dyadic window and adjust thresholds slightly to make
them continuity values.

The domains can be kept compactly away from the cylinder ends. The
Brownian-slab comparison of Proposition~\ref{conf:cm-comparison} gives a slab law
retaining positive probability for the witness, with $\overline W$ in
the slab interior, independent $J$, and the same field kernels. Restrict
the end-coordinate parameters only as needed for this room. Refine the
thresholds to continuity values in that law, and write $p_E>0$ for its
probability. Although $E,A,U,W$ may depend on $\eta$, the previously
chosen local recipe constants do not.

The exit safety implies that minimizing paths for these anchors are
attained in $K$, by compactness and the length topology. Such a geodesic
is almost surely unique up to parametrization. Here is an elementary
argument sufficient for these fixed anchors. Distinct geodesics can be
distinguished by a smooth nonnegative bump compactly supported in $W$,
missed by one and visited in its positive set by the other. Under addition
of $t$ times that bump, the distance is nondecreasing in $t$. At a
parameter with these two minimizers it is constant to the right, because
the avoiding minimizer is unchanged, and strictly smaller at every smaller
parameter, because the other minimizer has positive length in the bump's
positive set. There is at most one such parameter for a fixed field and
bump. A countable family of bumps distinguishes all distinct traces.
Local Weyl scaling, Brownian-slab shift equivalence, and Fubini then show
that the exceptional event has probability zero. The whole distance
family is obtained from the original metric by length reweighting, so
this parameter argument uses no additional metric construction or
confluence assertion.

\smallskip\noindent\emph{The input weight of successful visits.}\par
Integrate reference arclength along the middle half of this geodesic on
$E$. The aligned-shell estimate, including the successful shortcut
proposals and all the cutoffs above, gives order $N$ successful holes
around every point, with probability tending to one. For sufficiently
large $N$, the anchors lie outside every visited shell's enlarged disk,
and each disk has room inside $U$. Indeed the middle half is a fixed
positive reference distance from either anchor, while all spatial radii
tend uniformly to zero; use reference continuity, geodesicity, and the
compact safety condition.

Time spent in one hole along the middle half is at most a constant times
its length unit $w$. Replace the portion between first and last encounters
with the attachment circle by a path of cost at most $K_0w$ in its annular
buffer, and use geodesicity. Integrating with multiplicity therefore gives
the expected weight bound
\begin{equation}\label{metric:input-hit-weight}
 \E\!\left[\one_E\sum_{\text{successful hits}}w\right]
 \ge c_0 A p_E N
\end{equation}
for all sufficiently large $N$, with $c_0>0$ independent of $\eta$ and
$E$. Ordinary compact success events whose probabilities tend to one may
be imposed first. Their error is harmless even when $p_E$ is small,
since $E$ and its positive probability are fixed before $N\to\infty$.

\smallskip\noindent\emph{Domination for one transformed hit.}\par
For a successful hit, take a usable recipe from the finite list in
Lemma~\ref{metric:forcing-recipe}. One may count all usable bin and recipe
choices instead of prescribing a choice at an exact endpoint. For each
\emph{fixed} recipe and shell entry, its hit submeasure, together with the
retained inner and old exterior data, is bounded by $\Lambda$ times a
single common output law, where $\Lambda<\infty$ is independent of
$\eta,E,N$.

To prove this domination, first condition on the slab coordinate parameters
and fields, before imposing the hit. Keep the selected independently
proposed inner profile at $h+b$, and the old profile readings beyond a
strictly enlarged comparison disk. The inner preservation region has
$h+f=h+b$; the outer region has $h+f=h$. Source elimination and the
field-kernel comparison give exactly the corresponding two kernels of
the output law at $h'=h+f$. They have the conditional product law, by
\eqref{eq:original-15} after source elimination. Retain jointly all
countable regularized exterior readings with compact confinement in $U$
beyond a strictly containing comparison circle. Each finite collection
is separated from the inner protection and lies in the field-agreement
region, so compact exhaustion gives the whole collection. These kernel
identities hold for typical supported fields, including the deterministic
shifts in the finite list, and are independent of the end and source
choices.

Now apply Cameron--Martin to the fixed shift in the Brownian slab. Its
energy and pairing have the ordinary Dirichlet expression in the local
planar chart: the cylinder coordinate charge is harmonic on the small
support, and the conditional Brownian linear mean contributes no term.
The bounded-likelihood cutoff was part of shell success before selecting
anchors or a hit. Its bound is uniform in the slab endpoints and the
parameters restricted to have the required interior room. Imposing the
hit, choosing a covering bin, or imposing further restrictions only
decreases the resulting submeasure. Complete the output profile over its
retained readings using its usual conditional law; the same domination
then applies to all additional output variables used below. This is a
marginal comparison for one recipe and one entry. It makes no simultaneous
claim about transporting the auxiliary proposals for different entries
to one output realization. Gain and existence requirements can be
expressed by measurable countable tests with slightly weakened slack.

\smallskip\noindent\emph{Stability of the near-saturation event.}\par
The output satisfies the inequalities defining $E$ with errors tending
uniformly to zero over the eligible hits, after a negligible input
restriction if needed. First, anchor distances and exit-safety distances
change by $o(1)$: the shift is bounded and supported in one shrinking
interior ball. Cut a path from its first entrance to its last exit and
splice across that ball. Ordinary continuity of the old internal metric
and the bounded Weyl multiplier control both directions of this
comparison uniformly on an interior compact.

The certificate cost cannot decrease by a nonvanishing amount either.
Suppose an output certificate is cheaper. Mark its first entrance and
last exit on a circle farther out than the transferred-reading threshold,
still with the ball and buffers compactly inside $U$. If the certificate
stays exterior, its reading transfers directly. Otherwise approximate it
with losses tending to zero and extract the two approach points and
subbudgets. Its two exterior pieces have compact confinement exterior to
the threshold and are among the retained readings. The old $F_U$ can
use them and join the two approach points across the ball at cost $o(1)$,
by its triangle inequality and upper bound, together with the uniform
continuity of old internal costs on neighborhoods of the hit region.
Take infima over certificates, using countable intermediate protections
if necessary. Each candidate only needs its own compact confinement in
$U$; no common predetermined distance from $\partial U$ is required.
This proves
\[
 F'_U(p_0,p_1)\ge F_U(p_0,p_1)-o(1).
\]
This argument concerns open endpoint losses, and does not attach an
individual graph vertex to a small ball.

Old modulus and clearance cutoffs make all these errors deterministic
and uniformly vanishing over the shrinking hits, while preserving
\eqref{metric:input-hit-weight}. Hence the relaxed output event has
common-law probability $p_E+o(1)$, by continuity of the thresholds. Its
safety still gives minimizing paths in $K$ for large $N$.

\proofstage{A bound on the total output weight}
For each eligible transformed hit, the unique output geodesic has the
strict-gain interval and the in-disk traversal supplied by
Lemma~\ref{metric:forcing-recipe}. The per-label overlap tests of the
aligned-shell construction are unchanged by this one shift: its own
large outer circle average, the required neighbor averages, and the
coarser delay checks are all outside the support. Thus eligibility can
be defined using only necessary output properties: these unchanged
alignment tests, the relaxed event, and existence of the weakened gain
and traversal requirements. On every realization of the common output
law, eligible enlarged disks have bounded overlap for each label,
including the fixed finite recipe multiplicity.

Their marked intervals also have bounded overlap in geodesic time for
each label. All intervals containing a specified time lie in a span of
length at most $2B_0$, since each has length at most $B_0$. Each spends
at least \(a'u/2\) in its own enlarged disk. If the disks have overlap
at most \(M\) at this label and \(n_t\) intervals contain a given
time \(t\), summing these in-disk durations gives
\[
 n_t\,\frac{a'u}{2}\le 2M B_0,
 \qquad n_t\le\frac{4M B_0}{a'u}.
\]
The right-hand side is fixed because \(B_0/u\) is fixed. Partition the intervals into a bounded number of
families of disjoint intervals. Within a family, use the triangle
inequality for $F'_U$ to splice the improved pieces with the intervening
geodesic portions, charged at upper factor $C$. The total gain satisfies
\[
 \sum\delta_*u\le C L'-F'_U(p_0,p_1)
 \le O(\eta A)+o(1),
\]
where $L'=D_{h'}(p_0,p_1;W)$. The geodesic is compactly in $K\Subset U$,
so its subpaths are admissible for every such splice. Since
$w=2^{\xi m}u$ and $m$ is fixed, summing the bounded number of families,
then the labels in $[N,2N]$, gives
\begin{equation}\label{metric:output-hit-weight}
 \sum_{\text{eligible output hits}}w\le C_2\eta A N
\end{equation}
on the relaxed event, for all sufficiently large $N$. Here $C_2$ is
independent of $\eta,E,p_E$. Countable interval tests and slightly weaker
gains can be used throughout to make this bound measurable.

Apply the per-hit domination and integrate
\eqref{metric:output-hit-weight}. Together with
\eqref{metric:input-hit-weight} it yields
\[
 c_0A p_E N\le \Lambda C_2\eta A\bigl(p_E+o(1)\bigr)N.
\]
Choose $\eta>0$ sufficiently small first, obtaining its fixed event of
probability $p_E>0$, and then take $N$ sufficiently large. This is a
contradiction. Thus $c=C$, as claimed.
\end{proof}

\subsection{Quantile comparison and the scale bootstrap}

\begin{lemma}[Quantiles in a rigid extraction]\label{metric:rigid-quantiles}
For the rigid extractions of Proposition~\ref{metric:rigidity}, there
are deterministic constants $0<d_0<D_0<\infty$, independent of the
extraction and of each fixed integer offset $r$, such that
\begin{equation}\label{eq:original-17}
 d_0C\le\liminf\frac{A_{j-r}}{Tw_r}
 \le\limsup\frac{A_{j-r}}{Tw_r}\le D_0C.
\end{equation}
This includes forward offsets. In extractions with $T=A_j$ the common
factor $C$ lies in a universal compact subinterval of $(0,\infty)$.
There is also a fixed sufficiently small $d_1>0$ such that the lower
traversal test at threshold $d_1A_{j-r}$ strictly re-passes in these
extractions, with the required band slack.
\end{lemma}

\begin{proof}
Use the transported prototype law at offset $r$. Upper following of
finitely many reference winding instructions with strict clearance and
forced transverse stitching gives an upper probability cutoff for the
circuit cost at a fixed multiple of $C$. Lower comparison charges a
nontrivial compact subtraversal of every winding path, giving a strictly
positive lower probability cutoff at a fixed multiple of $C$. The
reference lengths in question have the same finite upper and positive
lower probability cutoffs at every offset. Together with the two
quantile inequalities \eqref{metric:quantile-tests}, these bounds give
\eqref{eq:original-17}. All margins are taken strict before returning
to the measurement band.

At a quantile-normalized label with finite upper comparison, the second
inequality in \eqref{metric:quantile-tests} already forces $C>0$ by
upper following, so Proposition~\ref{metric:rigidity} applies.
Putting $r=0$ and $T=A_j$ in \eqref{eq:original-17} gives
$1/D_0\le C\le1/d_0$. Finally, lower comparison by $C$ and the common
positive reference traversal cutoff give high-probability lower
success at a fixed small multiple of $CTw_r$. The upper bound in
\eqref{eq:original-17} makes $d_1A_{j-r}$ strictly smaller than this
threshold for sufficiently small fixed $d_1$. This gives the last
assertion with its probability and geometric slack.
\end{proof}

\begin{proposition}[Deterministic scale and local comparison]\label{metric:scale-bootstrap}
Define
\[
 \mathcal B(P)=A_{\lfloor\log_2P\rfloor}
\]
for large $P$, extending it positively for small $P$. For every
$\sigma\in(0,\xi)$ there is a finite deterministic $K_\sigma$ such that
\begin{equation}\label{eq:original-18}
 \frac{\mathcal B(P')}{\mathcal B(P)}
 \le K_\sigma\left(\frac{P'}P\right)^\sigma,
 \qquad 1\le P'\le P.
\end{equation}
Along every sequence of diverging perimeter units there is a further
joint extraction on which the lower and upper open-port comparison
factors in denominator units $\mathcal B(P)$ agree. Their common value
is deterministic and lies in a universal compact positive range.
These local almost-sure comparisons also hold on ordinary patches of
the unit-area sphere.
\end{proposition}

\begin{proof}
Fix $\sigma\in(0,\xi)$. At every sufficiently large index test the
following three assertions:
\begin{enumerate}
\item $M^{-1}<A_j/A_{j-1}<M$, for a fixed large $M$;
\item $A_{j-m_0}/A_j<2^{-\sigma m_0}$, for a fixed large integer $m_0$;
\item lower traversal success at threshold $d_1A_j$.
\end{enumerate}
Choose these constants so that rigid comparisons strictly re-pass all
three tests. Such a choice is possible by
Lemma~\ref{metric:rigid-quantiles}. In particular, choose $M$ larger
than the uniform bounds for the adjacent ratios. Then take $m_0$ so
large that the fixed multiplicative loss in
\eqref{eq:original-17}, times $2^{-\xi m_0}$, is strictly smaller
than $2^{-\sigma m_0}$. The lower test uses the fixed $d_1$ of that
lemma, decreased if necessary to retain strict slack.

Every fixed finite success run eventually holds about the record starts
of Lemma~\ref{metric:records}. Indeed, every extraction at such starts
has the upper ratio bound \eqref{metric:record-ratios} at all fixed
preceding offsets. Upper chaining gives a finite upper constant, and
quantile normalization makes it positive. Proposition~\ref{metric:rigidity}
and Equation~\eqref{eq:original-17} then apply at all fixed offsets,
including forward ones. A failure in a fixed finite run would contradict
their strict retests after further joint extraction.

Suppose nevertheless that there are later first exits after arbitrarily
long successful runs. Extract at the index $j$ immediately before such
a first failure. All fixed backward offsets, including zero, now pass
the three tests. The adjacent bounds permit a further extraction with
\begin{equation}\label{metric:ratio-sequence}
 r_i=\lim\frac{A_{j-i}}{A_jw_i}\in(0,\infty)
 \quad(i\ge0),\qquad r_0=1.
\end{equation}
We claim that these $r_i$ are uniformly comparable.

Fix a starting offset $i$. In $A_jw_i$ units, the farther tests at relative
offset $l$ give high upper and lower factors of order $r_{i+l}$ in
their exact-power units. If at least a fraction $1/10$ of the offsets
$l\in[N,2N]$ had $r_{i+l}\ge L_0r_i$ along arbitrarily large $N$,
lower chaining would give a lower factor of order $L_0r_i$ at $i$.
For a sufficiently large absolute $L_0$, this contradicts the winding
upper test there, whose scale is $r_i$. If instead the same fraction
had $r_{i+l}\le r_i/L_0$, upper chaining would give an upper factor
of order $r_i/L_0$, contradicting the second quantile inequality at $i$.
The chain constants and reference probability cutoffs are uniform, as
proved above. Thus, for each fixed $i$, a large majority of all
sufficiently long such intervals satisfy
\[
 r_i/L_0<r_{i+l}<L_0r_i.
\]
The large-majority sets for any two fixed starting indices overlap
after translating their intervals by the bounded difference of these
indices. A common member shows that their $r_i$'s differ by at most
$L_0^2$. Comparing with $r_0=1$ proves the claim.

At the index before failure, \eqref{metric:ratio-sequence} now gives an
upper ratio bound by a fixed multiple of $w_i$ at every fixed preceding
offset. Upper chaining again gives finite upper comparison; normalization
gives positivity; and rigidity follows. Equation~\eqref{eq:original-17}
and its lower traversal retest strictly re-pass all the tests at the
alleged failure index, a fixed forward offset. This is a contradiction.
Consequently all three tests hold eventually.

The same majority argument applies to any extraction at later indices:
the adjacent tests give positive finite fixed-offset limits, the upper
and lower tests make them uniformly comparable, and chaining followed
by rigidity gives equality of the comparison factors in $A_j$ units.
Lemma~\ref{metric:rigid-quantiles} bounds these factors uniformly above
and away from zero.

Iterate the second eventual test in blocks of length $m_0$ and use
the adjacent test for the fewer than $m_0$ remaining steps. After
enlarging a deterministic constant to include the finitely many small
indices, this gives
\[
 \frac{A_{j'}}{A_j}\le K\,2^{-\sigma(j-j')},\qquad j'\le j.
\]
Converting dyadic indices to $P',P$ costs only a fixed factor and proves
\eqref{eq:original-18}. In particular $\mathcal B(P)\to\infty$.

For arbitrary diverging perimeter units, their ratio to the lower
dyadic unit lies in $[1,2)$. Pass to a subsequence where it has a
positive limiting ratio, and apply the joint fixed-ratio unit conversion
of Lemma~\ref{kern:unit-conversion}. This transfers the same local statements,
with a deterministic common factor in a universal compact positive
range, after further extraction. Only local internal reference length
data occur in these assertions. The fixed-volume product rule
\eqref{eq:original-15}, retained by Proposition~\ref{kern:fixed-volume},
and the field matching and area-one transfer of
Lemmas~\ref{conf:field-kernel-locality} and~\ref{kern:area-one-transfer}
therefore transport them to ordinary patches of the unit-area sphere.
\end{proof}

\begin{remark}[Scope of the local conclusion]\label{metric:no-endpoint-control}
The upper statement follows complete finite-length reference paths with
arbitrary open loss. The lower statement applies to \emph{all} compactly
confined bounded-cost graph paths, and its positive factor also bounds
intrinsic reference distance in an open collar. Neither statement gives
a diameter bound for all vertices approaching a point, or the cost of
attaching an arbitrarily specified lattice endpoint. These additional
estimates are proved in Sections~\ref{sec:endpoints} and~\ref{sec:address}.
\end{remark}
